\setcounter{section}{0}

\Needspace{5\baselineskip}
\hypertarget{ux4e16ux754c-ux52a8ux4f5cux6a21ux578bux6982ux8ff0}{%
\section{世界-动作模型概述}\label{ux4e16ux754c-ux52a8ux4f5cux6a21ux578bux6982ux8ff0}}

\Needspace{5\baselineskip}
\hypertarget{ux4e00ux4e16ux754c-ux52a8ux4f5cux6a21ux578bux7684ux6982ux5ff5}{%
\subsection{一、世界-动作模型的概念}\label{ux4e00ux4e16ux754c-ux52a8ux4f5cux6a21ux578bux7684ux6982ux5ff5}}

世界-动作模型（World-Action Model，WAM）是世界模型范式在具身智能基础模型中的运用。世界-动作模型将状态预测和动作生成相结合，输入过去的状态（和已经采取的动作），用统一模型依次或同时预测下一状态和接下来的动作。其中，用统一模型是考虑到本质上都是在学习物理世界的规律，很多参数本就是重合的，且统一模型更符合端到端原则。

\Needspace{5\baselineskip}
\hypertarget{ux4e8cux4e16ux754c-ux52a8ux4f5cux6a21ux578bux7684ux8badux7ec3ux76eeux6807}{%
\subsection{二、世界-动作模型的训练目标}\label{ux4e8cux4e16ux754c-ux52a8ux4f5cux6a21ux578bux7684ux8badux7ec3ux76eeux6807}}

\Needspace{5\baselineskip}
WAM的损失函数一般为状态（视频像素或视觉特征）损失与动作损失加权求和：

\[
\mathcal{L}_{\mathrm{vid}}+\lambda_{\mathrm{act}}\mathcal{L}_{\mathrm{act}}
\]

\Needspace{5\baselineskip}
两项大多采取流匹配损失：

\Needspace{5\baselineskip}
对于目标变量 \(y\)（既可以是动作序列 \(a_{1:H}\)，也可以是未来视觉特征 \(z_{1:T}\)），算法采样高斯噪声 \(\varepsilon\sim\mathcal{N}(0,I)\) 和时间步 \(t\in(0,1)\)，构建插值样本：

\[
y_t=(1-t)y+t\varepsilon.
\]

\Needspace{5\baselineskip}
模型通过标准的流匹配损失预测速度场：

\[
\mathcal{L}_{\mathrm{FM}}(y)=\mathbb{E}_{y,\varepsilon,t}\left[\left\|f_\theta(y_t,t,o,l)-(\varepsilon-y)\right\|_2^2\right].
\]

\Needspace{5\baselineskip}
\hypertarget{ux4e09ux5177ux8eabux539fux751fux6a21ux578b}{%
\subsection{三、具身原生模型}\label{ux4e09ux5177ux8eabux539fux751fux6a21ux578b}}

\Needspace{5\baselineskip}
早期的WAM往往由视频生成模型后训练得到。但视频生成模型本来是为数字内容创作而生，搬到机器人身上，会带来三个问题：

\begin{enumerate}
\def\labelenumi{\arabic{enumi}.}
\item
  表征的VAE只为像素级重建优化，宝贵的潜空间大半让给了光影、背景纹理和物体表面的高光，而杯子正在向左移动这种真正跟控制有关的事，只占极小的一部分。它保住了视觉细节，丢掉的恰恰是物理和语义结构。且动作模块往往是一个外挂部分，世界状态和动作被留在两个对不齐的空间里。
\item
  速度高维的视频token叠上多步迭代的去噪过程，延时长，而真机操作要求策略在每个动作块的边界上读入新观测、随时修正。
\item
  互联网视频没有动作标签，而通用的视频预测目标并不教模型``动作如何改变世界''。
\end{enumerate}

故现在许多WAM转向了用具身数据原生训练。这种方式对算力和数据要求更高，但模型能力的上限也更高。

\FloatBarrier
\Needspace{5\baselineskip}
\hypertarget{ux53c2ux8003ux6587ux732e-158}{%
\subsection{参考文献}\label{ux53c2ux8003ux6587ux732e-158}}

\vspace{0.25em}
\begin{itemize}
\tightlist
\item
  Wang, S., Shi, J., Fu, Z., et al.~(2026). \href{https://arxiv.org/abs/2605.12090}{World Action Models: The Next Frontier in Embodied AI}. arXiv:2605.12090.
\item
  Lipman, Y., Chen, R. T. Q., Ben-Hamu, H., Nickel, M., \& Le, M. (2023). \href{https://arxiv.org/abs/2210.02747}{Flow Matching for Generative Modeling}. ICLR.
\end{itemize}

\clearpage
